\section*{第 \ref{ch:g12:quantum-world} 章——量子世界：光子与能级}

\begin{solution}{exo:g12:quantum-world:1}
$E = hc/\lambda = \num{1.99e-25}/\num{5.30e-7} = \qty{3.75e-19}{J} =
\qty{2.35}{eV}$. $N = \num{1.0e-3}/\num{3.75e-19} \approx \num{2.7e15}$
每秒光子数。
\end{solution}

\begin{solution}{exo:g12:quantum-world:2}
$E = h\nu = 6.63\times10^{-34} \times 10^{8} = \qty{6.63e-26}{J}
\approx \qty{4.1e-7}{eV}$——约为一个可见光光子能量的
五百万分之一。任何可探测的无线电信号都
包含每秒如此多的光子，因而其光流十分平滑。
\end{solution}

\begin{solution}{exo:g12:quantum-world:3}
$\nu_0 = W_0/h = \num{3.04e-19}/\num{6.63e-34} \approx
\qty{4.6e14}{Hz}$；$\lambda_0 = c/\nu_0 \approx \qty{654}{nm}$。取
\qty{633}{nm}，$E = 1240/633 = \qty{1.96}{eV} > \qty{1.9}{eV}$：能，
勉强，此时 $E_{k,\max} \approx \qty{0.06}{eV} \approx \qty{1e-20}{J}$。
\end{solution}

\begin{solution}{exo:g12:quantum-world:4}
$E_3 - E_2 = -1.51 + 3.40 = \qty{1.89}{eV}$;
$\lambda = 1240/1.89 \approx \qty{656}{nm}$：即红色谱线。
\end{solution}

\begin{solution}{exo:g12:quantum-world:5}
$\nu_0 \approx \qty{5.6e14}{Hz}$; $W_0 = h\nu_0 \approx
\qty{3.7e-19}{J} \approx \qty{2.3}{eV}$（即 $-W_0$ 截距）。
斜率即普朗克常量 $h$。
\end{solution}

\begin{solution}{exo:g12:quantum-world:6}
$E = 1240/400 = \qty{3.10}{eV}$; $E_{k,\max} = 3.10 - 2.3 =
\qty{0.80}{eV} = \qty{1.3e-19}{J}$; $U_s = E_{k,\max}/e =
\qty{0.80}{V}$.
\end{solution}

\begin{solution}{exo:g12:quantum-world:7}
光强加倍：电流加倍（光子数加倍），$U_s$
不变（每个光子能量相同）。频率升高：$U_s$ 增大，
$E_{k,\max} = h\nu - W_0$。波动图像将电子的能量
与光强相联系，而与频率无关——这恰恰是\emph{不}
被观察到的现象。
\end{solution}

\begin{solution}{exo:g12:quantum-world:8}
从 $n = 1$ 出发：\qty{13.6}{eV}，$\lambda = 1240/13.6 \approx
\qty{91}{nm}$，远紫外线。从 $n = 2$ 出发：\qty{3.40}{eV}，
$\lambda \approx \qty{365}{nm}$，近紫外线。
\end{solution}

\begin{solution}{exo:g12:quantum-world:9}
从基态出发，最小跃迁所需能量为 $E_2 - E_1 =
\qty{10.2}{eV}$（\qty{122}{nm}）：所有吸收谱线均
为紫外线，无一可见。在热的恒星大气中，碰撞使
部分原子处于 $n = 2$，从而能吸收可见的
巴耳末光子（\qtyrange{1.89}{3.40}{eV}）。
\end{solution}

\begin{solution}{exo:g12:quantum-world:10}
(a) $E_k = \qty{1.6e-17}{J}$, $p = \sqrt{2m_e E_k} =
\qty{5.4e-24}{kg.m/s}$, $\lambda = h/p \approx \qty{1.2e-10}{m}$ ---
原子尺度：晶体能使其发生衍射。(b) $p = \num{1.0e-4} \times 1.0 =
\qty{1.0e-4}{kg.m/s}$, $\lambda \approx \qty{6.6e-30}{m}$ --- $10^{20}$
倍（比原子小），这只苍蝇的波永远不可见。
\end{solution}

\begin{solution}{exo:g12:quantum-world:11}
一个光子：$\num{1.99e-25}/\num{5.05e-7} = \qty{3.94e-19}{J}$；九十个光子：
$\approx \qty{3.5e-17}{J}$。蚊子：$\tfrac12 \times \num{2.5e-6}
\times 0.50^2 = \qty{3.1e-7}{J}$——$10^{10}$ 倍。人眼
工作时距绝对极限仅相差百余个
光子，几乎就是一台光子计数器。
\end{solution}

\begin{solution}{exo:g12:quantum-world:12}
$\nu_1 = \qty{7.41e14}{Hz}$, $\nu_2 = \qty{5.49e14}{Hz}$;
$h = e\,\Delta U_s/\Delta\nu = \num{1.60e-19} \times 0.79 /
\num{1.91e14} \approx \qty{6.6e-34}{J.s}$. $W_0 = h\nu_2 - eU_{s,2} =
\num{3.63e-19} - \num{0.64e-19} = \qty{3.0e-19}{J} \approx
\qty{1.9}{eV}$：铯。
\end{solution}

\begin{solution}{exo:g12:quantum-world:13}
$1240/486 = \qty{2.55}{eV} = E_4 - E_2 = 3.40 - 0.85$：跃迁
$4 \to 2$. $1240/434 = \qty{2.86}{eV} = E_5 - E_2 = 3.40 - 0.54$:
跃迁 $5 \to 2$。
\end{solution}

\begin{solution}{exo:g12:quantum-world:14}
横截面 $\pi r^2 \approx \qty{3.1e-20}{m^2}$，收集的功率
$\qty{3.1e-26}{W}$; $t = \num{3.7e-19}/\num{3.1e-26} \approx
\qty{1.2e7}{s}$——约四个月，而观测到的时间不足一纳秒。
能量无法在波前上均匀铺开：它
集中地以颗粒形式传播。
\end{solution}

\begin{solution}{exo:g12:quantum-world:15}
$E_k = \qty{8.0e-16}{J}$, $p = \sqrt{2m_e E_k} =
\qty{3.8e-23}{kg.m/s}$, $\lambda = h/p \approx \qty{1.7e-11}{m} =
\qty{17}{pm}$——比 \qty{550}{nm} 约短 $\num{3e4}$ 倍。
衍射会使小于 $\sim\lambda$ 的细节变模糊，因此
将 $\lambda$ 缩小 $10^4$ 倍，分辨率便提高 $10^4$ 倍。
\end{solution}

\begin{solution}{pb:g12:quantum-world:1}
\textbf{1.} $E = hc/\lambda = \num{1.99e-25}/\num{5.50e-7} =
\qty{3.62e-19}{J} = \qty{2.26}{eV}$.

\textbf{2.} $N = \num{1.0e3}/\num{3.62e-19} \approx \num{2.8e21}$
个光子每秒。

\textbf{3.} $P = 0.20 \times \qty{1.0e3}{W} = \qty{200}{W}$;
$E = 200 \times 5.0 = \qty{1.0}{kW.h}$ ($\qty{3.6e6}{J}$).

\textbf{4.} 每秒 $10^{21}$ 个光子，统计
涨落极其微小，颗粒性被平均掉，
电流近乎完全平滑，正如快速倾泻的沙子流动如水。

\textbf{5.} 造成损伤只需一个能量为 \qtyrange{3}{4}{eV} 的光子：紫外
线光子合格，红外线光子（$\sim \qty{0.1}{eV}$）则永远不行——
数量再多也无法代替单个光子的威力。

\textbf{6.} $U_s$ 是恰好使最快的
光电子停下来的反向电压：$eU_s = E_{k,\max} = h\nu - W_0$。

\textbf{7.} 由 $\nu = c/\lambda$ 得：$8.22$、$7.41$、$6.88$ 与 $5.49$
$\times\,10^{14}\,\unit{Hz}$.

\textbf{8.} $eU_s = h\nu - W_0$ 是 $\nu$ 的一次函数，斜率为 $h$。
相邻两段斜率：$(0.94 - 0.37)/1.39 = 0.41$ 与
$(1.50 - 0.94)/1.34 = 0.42$ \unit{eV}，每 $10^{14}\,\unit{Hz}$——
三点共线。

\textbf{9.} $h = e(U_{s,1}-U_{s,4})/(\nu_1-\nu_4) = \num{1.60e-19}
\times 1.13/\num{2.72e14} = \qty{6.64e-34}{J.s}$.

\textbf{10.} 在 \qty{546}{nm} 时：$W_0 = h\nu - eU_s = \num{3.64e-19} -
\num{0.59e-19} = \qty{3.05e-19}{J} = \qty{1.90}{eV}$;
$\nu_0 = W_0/h = \qty{4.60e14}{Hz}$; $\lambda_0 \approx \qty{653}{nm}$.

\textbf{11.} 只有 $\lambda < \qty{653}{nm}$ 的光才能打出光电子：紫至红
可见光中从紫到红的部分以及紫外线。整个红外一侧——约占太阳能量的一半——
对该阴极毫无用处。

\textbf{12.} $(6.64 - 6.63)/6.63 \approx 0.2\%$——千分之二的
普适常量的一次物理实验室测量。

\textbf{13.} $E_2 = \qty{-3.40}{eV}$, $E_3 = \qty{-1.51}{eV}$,
$E_4 = \qty{-0.85}{eV}$, $E_5 = \qty{-0.54}{eV}$.

\textbf{14.} $3 \to 2$: \qty{1.89}{eV}; $4 \to 2$: \qty{2.55}{eV};
$5 \to 2$: \qty{2.86}{eV}.

\textbf{15.} $\lambda = 1240/E$：\qty{656}{nm}（红），\qty{486}{nm}
（蓝绿），\qty{434}{nm}（紫）。

\textbf{16.} 原子的能量是量子化的，因此只能发射
分立的能量差 $E_p - E_n$：能级阶梯产生的是谱线，
而不是连续的彩虹。

\textbf{17.} 否：\qty{500}{nm} 对应 \qty{2.48}{eV}，与任何能级差都不吻合
（$E_2$ 向上给出 $1.89$、$2.55$、$2.86$、
\qty{3.40}{eV}；从 $E_1$ 出发则至少需 \qty{10.2}{eV}）
——气体在该波段是透明的。

\textbf{18.} 由 $eU = \tfrac12 m_e v^2$：$v = \sqrt{2 \times
\num{1.60e-19} \times \num{1.0e4}/\num{9.11e-31}} \approx
\qty{5.9e7}{m/s} \approx 0.20\,c$——经典处理可以接受。

\textbf{19.} $p = m_e v = \qty{5.4e-23}{kg.m/s}$;
$\lambda = h/p = \qty{1.2e-11}{m} = \qty{12}{pm}$.

\textbf{20.} 增益 $= \num{5.50e-7}/\num{1.23e-11} \approx \num{4.5e4}$。
周末问题中的两个数值：$h = \qty{6.64e-34}{J.s}$，测量精确到
$0.2\%$ 以内，借助滤光片和电压表；另有一台显微镜，
分辨率高出四万五千倍，因为电子具有波动性。
\end{solution}
